% Originally: content/week-01.tex, \section{Young's inequality}
% Authored for these notes around the statement of problem 1.4; no tutor's notes cover it.
\section{Young's inequality}
\label{sec:youngs_inequality}

The following inequality is not part of the lecture's main line of development, but it is the
cleanest route to comparing different norms on $\mathbb{R}^n$ (see \cref{sec:cauchy_schwarz}) and
to the Hölder-type estimates that appear whenever two $L^p$-type quantities are played off against
each other. It is worth proving once and citing from then on.

\begin{proposition}[Young's inequality]
\label{prop:youngs_inequality}
% Source: exercises/Ex1_Analysis2_eng.pdf, Problem 1.4
Let $a, b > 0$ and let $p, q > 1$ satisfy $\tfrac{1}{p} + \tfrac{1}{q} = 1$. Then
\[
  ab \leq \frac{a^p}{p} + \frac{b^q}{q}.
\]
\end{proposition}

\begin{remark}
The case $p = q = 2$ recovers the familiar inequality $ab \leq \tfrac{a^2}{2} + \tfrac{b^2}{2}$,
which is just a rewriting of $(a-b)^2 \geq 0$.
\end{remark}

\begin{proof}
Fix conjugate exponents $p, q > 1$, and define $g : (0,\infty) \to \mathbb{R}$ by
$g(t) := \tfrac{t}{p} + \tfrac{1}{q} - t^{1/p}$. We first show that $g \geq 0$, with equality only
at $t = 1$. Differentiating,
\[
  g'(t) = \frac{1}{p}\left(1 - t^{-(p-1)/p}\right).
\]
Since $p > 1$, the exponent $-(p-1)/p$ is negative, so $t^{-(p-1)/p} > 1$ for $t \in (0,1)$ and
$t^{-(p-1)/p} < 1$ for $t > 1$. Hence $g' < 0$ on $(0,1)$ and $g' > 0$ on $(1,\infty)$, so $g$
attains a global minimum at $t = 1$. As $g(1) = \tfrac{1}{p} + \tfrac{1}{q} - 1 = 0$, we obtain
\begin{equation}
\label{eq:young_scalar_key}
  t^{1/p} \leq \frac{t}{p} + \frac{1}{q} \qquad \text{for all } t > 0.
\end{equation}
Now fix $a, b > 0$ and apply \eqref{eq:young_scalar_key} with $t := a^p / b^q$. Multiplying both
sides by $b^q > 0$ gives
\[
  \left(\frac{a^p}{b^q}\right)^{1/p} b^q \ \leq\ \frac{a^p}{p} + \frac{b^q}{q}.
\]
It remains to identify the left-hand side, and this is where the conjugacy of $p$ and $q$ enters.
Using $1 - \tfrac{1}{p} = \tfrac{1}{q}$,
\[
  \left(\frac{a^p}{b^q}\right)^{1/p} b^q
    \ =\ a\,b^{-q/p}\,b^{q}
    \ =\ a\,b^{\,q\left(1 - \frac{1}{p}\right)}
    \ =\ a\,b^{\,q \cdot \frac{1}{q}}
    \ =\ ab,
\]
which is exactly the asserted inequality.
\end{proof}

% Generator: Claude Opus 5 (High)
\begin{remark}[From Young to Hölder]
\label{rem:young_to_hoelder}
The exercise sheet also asks for the vector version
\begin{equation}
\label{eq:young_vector}
  \langle a,b\rangle \leq \frac{\|a\|_p^p}{p} + \frac{\|b\|_q^q}{q}, \qquad a,b \in \mathbb{R}^n,
\end{equation}
where $\|x\|_p := \left(\sum_i |x_i|^p\right)^{1/p}$. This is \emph{not} yet Hölder's inequality:
it is the scalar case applied coordinatewise and summed (see the solution of \cref{ex:1.4}
below), and the right-hand side is a sum of $p$-th and $q$-th powers rather than the product
$\|a\|_p\|b\|_q$. The step from one to the other is a normalisation, and it is worth seeing
once, because the same trick recurs whenever an inequality is homogeneous.

Assume $a,b \neq 0$ and apply \eqref{eq:young_vector} to the rescaled vectors
$a/\|a\|_p$ and $b/\|b\|_q$, each of which has $p$- resp.\ $q$-norm equal to $1$:
\[
  \frac{\langle a,b\rangle}{\|a\|_p\,\|b\|_q} \ \leq\ \frac{1}{p} + \frac{1}{q} \ =\ 1,
\]
which is exactly \newterm{Hölder's inequality} \germanfor{Höldersche Ungleichung}
$\langle a,b\rangle \leq \|a\|_p\|b\|_q$. The conjugacy condition
$\tfrac1p+\tfrac1q=1$, which so far looked like an arbitrary bookkeeping device, is what makes
the two constants add up to exactly $1$ here.

This is also a first hint of why $p$-norms other than $p=2$ are worth naming:
\Cref{sec:cauchy_schwarz} only treats $p = q = 2$, where Hölder's inequality is
Cauchy--Schwarz, but the same machinery bounds $\langle a, b\rangle$ for every conjugate pair
$(p,q)$.
\end{remark}

% Source: exercises/Ex1_Analysis2_eng.pdf, p. 2
\begin{exercise}[1.4 --- Young's inequality]
\label{ex:1.4}
Let $a,b > 0$, and let $p,q>1$ with $\tfrac{1}{p}+\tfrac{1}{q}=1$. Prove that
$ab \leq \tfrac{a^p}{p} + \tfrac{b^q}{q}$.
\begin{enumerate}[label=\textbf{\arabic*.}]
  \item Prove that for all $t>0$, $t^{1/p} \leq \tfrac{t}{p} + \tfrac{1}{q}$.
  \item Multiply the inequality by $b^q$ and choose an appropriate value of $t$ to conclude.
\end{enumerate}
\textbf{A bonus:} show that $\langle a,b\rangle \leq \tfrac{\|a\|_p^p}{p} + \tfrac{\|b\|_q^q}{q}$
for all $a,b \in \mathbb{R}^n$.
\end{exercise}
